\lampiransection{Dokumentasi Kinerja Nasi Gerilya dan Kompetitor pada GrabFood}
\label{app:dokumentasi-indeks-kinerja}

Lampiran ini memuat dokumentasi kinerja untuk penelitian \textit{Strategi Pemasaran Rumah Makan Nasi Gerilya pada Platform GrabFood untuk Meningkatkan Penjualan Berdasarkan Analisis SWOT}. Data yang disajikan mencakup jumlah penjualan bulanan dalam satuan bungkus, indeks \textit{Average Order Value} (AOV), indeks kunjungan harian, dan indeks kompetitor pada platform GrabFood. Jumlah penjualan bulanan disajikan dalam satuan bungkus sesuai grafik penjualan per bulan, sedangkan indeks AOV, indeks kunjungan harian, dan indeks kompetitor disajikan sebagai data turunan untuk membaca arah performa usaha dan posisi pembanding \textit{merchant}.

Data indeks digunakan karena sebagian data performa bersifat sensitif dan tidak seluruhnya dapat disajikan dalam nominal absolut. Indeks dihitung dengan menetapkan periode dasar sebesar 100,0, kemudian membandingkan nilai periode berikutnya terhadap periode dasar tersebut. Rumus umum yang digunakan adalah sebagai berikut.

\[
    \mbox{Indeks periode } t = \frac{\mbox{Nilai periode } t}{\mbox{Nilai periode dasar}} \times 100
\]

Pada Tabel~\ref{tab:indeks_aov_grabfood}, periode dasar adalah Juli 2025. Pada Tabel~\ref{tab:indeks_kunjungan_grabfood}, periode dasar adalah September 2025. Pada Tabel~\ref{tab:indeks_kompetitor_grabfood}, periode dasar adalah Juli 2025 untuk setiap \textit{merchant}. Perubahan persentase dihitung dari selisih indeks antarperiode dibandingkan indeks periode awal.

\begin{table}[!htbp]
    \centering
    \caption{Jumlah Penjualan Bulanan Nasi Gerilya}
    \label{tab:jumlah_penjualan_bungkus_grabfood}
    \begin{tabular}{|l|c|c|}
        \toprule
        \textbf{Periode} & \textbf{Jumlah Penjualan (Bungkus)} & \textbf{Perubahan} \\ \hline
        \midrule
        Des-24           & 6.702                               & --                 \\ \hline
        Jan-25           & 8.740                               & +30,4\%            \\ \hline
        Feb-25           & 8.911                               & +2,0\%             \\ \hline
        Mar-25           & 16.590                              & +86,2\%            \\ \hline
        Apr-25           & 15.129                              & -8,8\%             \\ \hline
        Mei-25           & 14.580                              & -3,6\%             \\ \hline
        Jun-25           & 13.467                              & -7,6\%             \\ \hline
        Jul-25           & 14.633                              & +8,7\%             \\ \hline
        Ags-25           & 14.259                              & -2,6\%             \\ \hline
        Sep-25           & 14.034                              & -1,6\%             \\ \hline
        Okt-25           & 13.009                              & -7,3\%             \\ \hline
        Nov-25           & 12.900                              & -0,8\%             \\ \hline
        \bottomrule
    \end{tabular}
    \tablesource{Dokumentasi grafik penjualan bulanan Nasi Gerilya; diolah peneliti (2026).}
\end{table}

\begin{table}[!htbp]
    \centering
    \caption{Indeks \textit{Average Order Value} (AOV)}
    \label{tab:indeks_aov_grabfood}
    \begin{tabular}{|l|c|c|}
        \toprule
        \textbf{Periode} & \textbf{Indeks} & \textbf{Perubahan} \\ \hline
        \midrule
        Jul-25           & 100,0           & --                 \\ \hline
        Ags-25           & 102,0           & +2,0\%             \\ \hline
        Sep-25           & 107,9           & +5,8\%             \\ \hline
        Okt-25           & 110,0           & +1,9\%             \\ \hline
        Nov-25           & 114,5           & +4,1\%             \\ \hline
        \bottomrule
    \end{tabular}
    \tablesource{Dokumentasi internal Nasi Gerilya; diolah peneliti (2026).}
\end{table}

\begin{table}[!htbp]
    \centering
    \caption{Indeks Kunjungan Harian}
    \label{tab:indeks_kunjungan_grabfood}
    \begin{tabular}{|l|c|c|}
        \toprule
        \textbf{Periode} & \textbf{Indeks} & \textbf{Perubahan} \\ \hline
        \midrule
        Sep-25           & 100,0           & --                 \\ \hline
        Okt-25           & 116,5           & +16,5\%            \\ \hline
        Nov-25           & 130,7           & +12,2\%            \\ \hline
        \bottomrule
    \end{tabular}
    \tablesource{Dokumentasi internal Nasi Gerilya; diolah peneliti (2026).}
\end{table}

\begin{table}[!htbp]
    \centering
    \small
    \caption{Indeks Kompetitor Nasi Padang Juli--Oktober 2025}
    \label{tab:indeks_kompetitor_grabfood}
    \begin{tabularx}{\textwidth}{|Y|>{\centering\arraybackslash}p{1.15cm}|>{\centering\arraybackslash}p{1.15cm}|>{\centering\arraybackslash}p{1.15cm}|>{\centering\arraybackslash}p{1.15cm}|>{\centering\arraybackslash}p{1.75cm}|}
        \toprule
        \textbf{\textit{Merchant}} & \textbf{Jul} & \textbf{Ags} & \textbf{Sep} & \textbf{Okt} & \textbf{\shortstack{$\Delta$ Jul--Okt}} \\ \hline
        \midrule
        Koto Gadang                & 100,0        & 105,1        & 86,2         & 94,2         & -5,8\%                                  \\ \hline
        Pondok Krakatau            & 100,0        & 110,1        & 89,0         & 107,5        & +7,5\%                                  \\ \hline
        Uda Sayang                 & 100,0        & 94,1         & 89,0         & 88,9         & -11,1\%                                 \\ \hline
        Dendeng Batokok            & 100,0        & 104,1        & 85,2         & 95,1         & -4,9\%                                  \\ \hline
        Istana Krakatau            & 100,0        & 119,2        & 108,4        & 119,3        & +19,3\%                                 \\ \hline
        Sederhana                  & 100,0        & 97,5         & 94,9         & 98,6         & -1,4\%                                  \\ \hline
        Gumarang Jaya              & 100,0        & 116,1        & 102,6        & 114,1        & +14,1\%                                 \\ \hline
        Nasi Gerilya               & 100,0        & 92,6         & 87,8         & 70,4         & -29,6\%                                 \\ \hline
        Garuda                     & 100,0        & 95,5         & 90,1         & 93,9         & -6,1\%                                  \\ \hline
        Pondok Gurih               & 100,0        & 102,9        & 97,0         & 101,5        & +1,5\%                                  \\ \hline
        \bottomrule
    \end{tabularx}
    \tablesource{Dokumentasi internal dan pembanding \textit{merchant}; diolah peneliti (2026).}
\end{table}

\begin{table}[!htbp]
    \centering
    \scriptsize
    \caption{Lokasi Nasi Gerilya dan Kompetitor Pembanding Berdasarkan Google Maps}
    \label{tab:lokasi_kompetitor_pembanding}
    \begin{tabularx}{0.96\textwidth}{|>{\RaggedRight\arraybackslash}p{2.6cm}|Y|>{\centering\arraybackslash}p{2.6cm}|}
        \toprule
        \textbf{\textit{Merchant}} & \textbf{Alamat pada Google Maps} & \textbf{\shortstack{Jarak\\Garis Lurus\\dari Nasi Gerilya}} \\ \hline
        \midrule
        Nasi Gerilya               & Jl. KL. Yos Sudarso No.111--111a, Glugur Kota, Kec. Medan Barat, Kota Medan, Sumatera Utara 20238       & --        \\ \hline
        Koto Gadang                & Jl. Gunung Krakatau No.12A--12B, Glugur Darat II, Kec. Medan Timur, Kota Medan, Sumatera Utara 20238    & 1,04 km   \\ \hline
        Pondok Krakatau            & Jl. Gunung Krakatau No.53BC, Glugur Darat I, Kec. Medan Timur, Kota Medan, Sumatera Utara 20238         & 0,93 km   \\ \hline
        Uda Sayang                 & Jl. Gunung Krakatau No.11, Pulo Brayan Darat II, Kec. Medan Timur, Kota Medan, Sumatera Utara 20239     & 1,68 km   \\ \hline
        Dendeng Batokok            & Jl. Kapten Muchtar Basri No.12, Glugur Darat II, Kec. Medan Timur, Kota Medan, Sumatera Utara 20238     & 0,45 km   \\ \hline
        Istana Krakatau            & Jl. Gunung Krakatau No.9, Pulo Brayan Darat II, Kec. Medan Timur, Kota Medan, Sumatera Utara 20239      & 1,62 km   \\ \hline
        Sederhana                  & Jl. Krakatau Ujung No.9G, Pulo Brayan Bengkel Baru, Kec. Medan Timur, Kota Medan, Sumatera Utara 20237  & 1,99 km   \\ \hline
        Gumarang Jaya              & Jl. Brigjend Katamso No.27--29, A U R, Kec. Medan Maimun, Kota Medan, Sumatera Utara 20212              & 4,02 km   \\ \hline
        Garuda                     & Jl. H. Adam Malik D/H Glugur ByPass No.14, Silalas, Kec. Medan Barat, Kota Medan, Sumatera Utara 20114  & 2,20 km   \\ \hline
        Pondok Gurih               & Jl. Brigjend Katamso No.33A, A U R, Kec. Medan Maimun, Kota Medan, Sumatera Utara 20159                 & 4,06 km   \\ \hline
        \bottomrule
    \end{tabularx}
    \tablesource[0.96\textwidth]{Observasi Google Maps pada 10 Juni 2026; jarak garis lurus dihitung peneliti dari koordinat Google Maps.}
\end{table}

Data pada Tabel~\ref{tab:jumlah_penjualan_bungkus_grabfood}, \ref{tab:indeks_aov_grabfood}, \ref{tab:indeks_kunjungan_grabfood}, \ref{tab:indeks_kompetitor_grabfood}, dan \ref{tab:lokasi_kompetitor_pembanding} digunakan sebagai dasar pembacaan arah performa usaha, posisi pembanding, dan kedekatan lokasi kompetitor pada platform GrabFood.
